\documentclass[11pt,a4paper]{article}

% Essential packages
\usepackage[utf8]{inputenc}
\usepackage[T1]{fontenc}
\usepackage[margin=1in]{geometry}
\usepackage{amsmath,amsfonts,amssymb}
\usepackage{graphicx}
\usepackage{booktabs}
\usepackage{url}
\usepackage{hyperref}
\usepackage{natbib}
\usepackage{setspace}


% Additional useful packages
\usepackage{fancyhdr}
\usepackage{subcaption}
\usepackage{tikz}
\usepackage{algorithm}
\usepackage{algorithmic}
\usepackage{listings}
\usepackage{xcolor}

% Configure hyperref
\hypersetup{
    colorlinks=true,
    linkcolor=blue,
    filecolor=magenta,      
    urlcolor=cyan,
    citecolor=red,
    pdftitle={Your Paper Title},
    pdfauthor={Your Name},
}

% Set line spacing
\onehalfspacing

% Configure headers and footers
\pagestyle{fancy}
\fancyhf{}
\rhead{\thepage}
\lhead{\textit{Location and Career Search}}
\renewcommand{\headrulewidth}{0.4pt}

% Title page information
\title{Comfort or Cost: \\
The Motivation for Migration After 40}

\author{
    Tate Mason\thanks{Tate.Mason@uga.edu} \\
    \textit{John Munro Godfrey Sr. Department of Economics} \\
    \textit{The University of Georgia} \\
    \textit{Athens, Georgia} \\
    \\
}

\date{\today}

% Abstract environment
\newenvironment{abstract}%
{\cleardoublepage\null \vfill\begin{center}%
\bfseries \abstractname \end{center}}%
{\vfill\null}

\begin{document}

% Title page
\maketitle
\thispagestyle{empty}

% Abstract
\begin{abstract}
\noindent % Write a concise summary (150-250 words) covering: research problem, methodology, key findings, and significance. State your main contributions clearly.

This work investigates adds to a broad literature on location search as well as the literature examining differences between workers and retirees.
Looking at individuals in the middle-late stage of their working careers, I seek to understand how motivations for relocating may differ from retirees.
Using data from the International Public Use Microdata Series (IPUMS) and the American Community Survey (ACS), I analyze patterns of migration among workers aged 40-80,
focusing on the differing motivations and outcomes for movers in this age group. Using structural modeling, I attempt to simulate the decision-making process of these individuals,
incorporating non-career factors like housing conditions and location-specific amentities. The findings suggest that those still working are likely to move to reduce cost of living
while also hoping to maintain or improve their career prospects, whereas retirees prioritize quality of life factors. This research contributes to a better understanding of the complex interplay
between career and non-career factors in relocation decisions among older adults, with implications for policymakers and urban planners aiming to attract and retain this demographic.

\vspace{0.3cm}
\noindent \textbf{Keywords:} % List 3-6 keywords relevant to your research
Location search, retirement decisions, career outcomes.
\end{abstract}

\newpage
\setcounter{page}{1}

% Main content
\section{Introduction}
\label{sec:introduction}

Much of the existing literature on location search focuses on young adults and their career-building moves. However, this leaves a gap in understanding how older adults,
particularly those in the later stages of their careers or approaching retirement, make relocation decisions. This paper aims to fill that gap by examining the motivations
and outcomes of location search among individuals aged 40-80. This demographic is particularly interesting as they may have different priorities compared to younger movers,
balancing career considerations with quality of life factors as they near retirement. Understanding these dynamics is crucial for policymakers and urban planners who seek to attract and retain older adults,
a demographic that is growing in many regions due to aging populations. 

In this paper, I utilize data from the International Public Use Microdata Series (IPUMS) and the American Community Survey (ACS) to analyze migration patterns among older adults. I look to identify the non-career
factors that influence their relocation decisions, such as housing conditions, cost of living, and access to amenities. By employing structural modeling techniques, I aim to simulate the decision-making process of
these individuals, providing insights into how they weigh various factors when considering a move.

This research contributes to the broader literature on migration and career search by highlighting the unique considerations of older adults. The findings have implications for urban development strategies,
as attracting this demographic can have significant economic and social benefits for communities. Further, contribution to the retirement literature is made by exploring how relocation decisions differ once
an agent exits the workforce. By creating an index of location-specific amenity factors that influence location choice, this work provides a foundation for future research into the complex interplay between 
career and non-career motivations in relocation decisions. 

Section II of this paper reviews related work in the field. Section III details the methodology employed in this study.
Section IV presents the analysis procedure, including datasets, modeling, calibration, results, and analysis. Section V discusses
the findings and potential future work. Finally, Section VI concludes the paper.

\section{Related Work}
\label{sec:related_work}

% Comprehensively review relevant prior work. Organize thematically, not chronologically.
% Show how your work builds on, differs from, or improves upon existing approaches.

\subsection{Previous Approaches}
\label{subsec:previous_approaches}

% Survey existing methods and approaches relevant to your problem.
% Group similar approaches and explain their key ideas, strengths, and weaknesses.

\subsection{Limitations of Existing Work}
\label{subsec:limitations}

% Identify specific gaps, limitations, or open problems in existing approaches.
% This motivates why your work is needed and positions your contributions.

\section{Methodology}
\label{sec:methodology}

% Describe your approach in sufficient detail for reproducibility.
% Include mathematical formulations, algorithms, and theoretical analysis as needed.

\subsection{Model Overview}
\label{subsec:formulation}

There are two types of agents in the model: workers and retirees. Workers will supply labor inelastically, for which they will earn a wage depending on the state they are in, $w(z)$. Each
period, they will retire with probability $1-x$. Retirees will receive pension income, $p$, and will die with probability $1-y$ each period, being replaced by a new worker. Both types of 
agents will maximize lifetime utility over consumption, $c$, and amenities, $m(z)$, in their location, as well as experience disutility from moving, $\gamma(z)$. 

\subsection{Utility Function}
\label{subsec:utility}
The utility functions for workers and retirees are defined as follows:
\begin{equation}
  U^i(c, m(z);t) = \frac{c^{1-\sigma}}{1-\sigma} - \gamma\cdot\{z'\ne z\} \ \forall \ i \in \{w, r\}
\end{equation}

where $c$ is consumption, $m(z)$ represents amenities in location $z$, and $\gamma$ is the moving cost parameter. The worker and retiree share common preferences over consumption and amenties, but
differ in their income sources and decision-making horizons. The constraint for workers is given by:

\begin{equation}
c_t = w(z_t) - \tau(z_t) - \mathbb{E}[c(h(z))|\delta]
\end{equation}

where $\delta$ represents prior assumptions about housing costs in location $z'$, $w(z_t)$ is the wage in location $z_t$, $\tau(z_t)$ are taxes, and $c(h(z))$ is the housing cost. $\delta$ can take
one of three values: optimistic, neutral, or pessimistic, representing different beliefs about housing costs in state $z'$. These beliefs will be updated each period, and the agents
will not hold onto prior beliefs. Thus, one can think of $\delta$ as a sufficient statistic for the agent's belief about housing costs in each state.

For retirees, the budget constraint is:

\begin{equation}
  c_t = p\times w_t - \mathbb{E}[c(hc(z))|\delta]
\end{equation}

where $p\times w_t$ is the pension replacement rate. The health care cost, $c(hc(z))$, is also influenced by the agent's belief $\delta$ about costs in location $z'$.
Similarly, $\delta$ can take one of three values: optimistic, neutral, or pessimistic, representing different beliefs about health care costs in location $z'$. These beliefs will also 
be updated each period and will not be held onto from prior periods.

\subsection{Decision Problem: Workers}
\label{subsec:worker_decision}
The value function for workers is defined as:
\begin{equation}
  V^w(\delta, z_t, t) = \max\{V_w^M(\delta, z_t, t), V_w^S(\delta, z_t, t)\}
\end{equation}
where $V_w^M$ is the value of moving and $V_w^S$ is the value of staying. The value of staying is given by:
\begin{equation}
  V^S_w(\delta, z_t, t) = \max_{c_t} \{U(c_t) + x\beta V^w(\delta, z_t, t+1) + (1-x)\beta V^r(\delta,z_t, t+1)\}
\end{equation}
The value of moving is given by:
\begin{equation}
  V^M_w(\delta, z_t; t) = \max_{z', c_t} \{U(c_t, m(z')) + x\beta \mathbb{E}[V^w(\delta', z'; t+1)|\delta] + (1-x)\beta\mathbb{E}[V^r(\delta', z'; t+1)|\delta]\}
\end{equation}

Here, $\delta'$ represents the updated belief about housing costs in the new location $z'$. If a worker stays, they retain their current belief $\delta$; if they move, they form a new belief $\delta'$
based on the observed costs in the new location and discard priod beliefs $\delta$. 

\subsection{Decision Problem: Retirees}
\label{subsec:retiree_decision}
The value function for retirees is defined as:
\begin{equation}
  V^r(\delta, z_t; t) = \max\{V_r^M(\delta, z_t; t), V_r^S(z_t; t)\}
\end{equation}
where $V_r^M$ is the value of moving and $V_r^S$ is the value of staying. The value of staying is given by:
\begin{equation}
  V^S_r(z_t; t) = \max_{c_t} \{U(c_t, m(z_t)) + y\beta V^r(\delta, z_t; t+1)\}
\end{equation}
The value of moving is given by:
\begin{equation}
  V^M_r(\delta, z_t; t) = \max_{z', c_t} \{U(c_t, m(z')) + y\beta \mathbb{E}[V^r(\delta', z'; t+1)|\delta]\}
\end{equation}
Similar to workers, retirees update their beliefs about health care costs in the new location $z'$ upon moving, forming a new belief $\delta'$ and discarding prior beliefs $\delta$.

\subsection{Proposed Approach}
\label{subsec:approach}

To solve the model, I will employ dynamic programming techniques to compute the value functions for both workers and retirees. The state space will include the current location $z_t$,
the belief state $\delta$, and the time period $t$. I will discretize the state space to make the problem computationally tractable. I will use value function iteration to solve for the optimal policies for moving or staying
for both types of agents. The model will be calibrated using observed migration patterns from the IPUMS and ACS datasets, adjusting parameters such as moving costs, wage functions,
and pension replacement rates to match empirical data from FRED. 


\subsection{Datasets}
\label{subsec:datasets}

I will utilize data from the International Public Use Microdata Series (IPUMS) and the American Community Survey (ACS) to analyze migration patterns among individuals aged 40-80. These datasets provide rich demographic and
socioeconomic information, allowing for a comprehensive analysis of relocation decisions in this age group. I will extract relevant variables such as age, income, housing conditions, and migration history to inform the model calibration and validation.
Further, using data from FRED, I will calibrate macroeconomic parameters such as wage growth, pension replacement rates, and cost of living indices to ensure the model accurately reflects real-world conditions. I have also created an index of location-specific
amenity factors that influence location choice, which will be incorporated into the model to capture non-career motivations for relocation. These amenity factors include climate, healthcare quality, recreational opportunities, and cultural amenities. These
factors are sourced from...

\subsection{Summary Statistics}
\label{subsec:summary_statistics}


\section{Calibration and Results}
\label{sec:results_analysis}

\subsection{Calibration}
\label{subsec:calibration}

To calibrate the model, I will use a combination of method of moments and maximum likelihood estimation. The calibration process will involve adjusting model parameters to match key empirical moments observed in the IPUMS and ACS datasets.
These moments include migration rates, moving decisions based on age and income, and the distribution of movers across different locations. I will also calibrate parameters related to moving costs, wage functions, and pension replacement rates
using data from FRED. The calibration will be performed iteratively, refining parameters until the model closely aligns with observed data.

I will validate the calibrated model by comparing its predictions to out-of-sample data, ensuring that it accurately captures the dynamics of relocation decisions among older adults.

\subsection{Results}
\label{subsec:results}

% Present results clearly and objectively:
% - Use tables and figures effectively
% - Report statistical significance where appropriate
% - Include error bars or confidence intervals
% - Compare against all relevant baselines
% - Report both quantitative metrics and qualitative observations

\subsection{Analysis}
\label{subsec:analysis}

% Provide deeper analysis of your results:
% - Statistical significance tests
% - Ablation studies showing contribution of different components
% - Error analysis and failure cases
% - Computational complexity and runtime analysis
% - Sensitivity analysis for key parameters

\section{Discussion}
\label{sec:discussion}

% Interpret your findings, discuss broader implications, and acknowledge limitations honestly.

\subsection{Interpretation of Results}
\label{subsec:interpretation}

% Explain what your results mean in the context of your research questions:
% - How do findings relate to your hypotheses?
% - What insights do the results provide?
% - How do they compare to previous work?
% - What are the practical implications?

\subsection{Future Work}
\label{subsec:future_work}

% Suggest concrete directions for future research:
% - How could your approach be extended or improved?
% - What new research questions arise from your work?
% - How could limitations be addressed?
% - What are the next logical steps?

\section{Conclusion}
\label{sec:conclusion}

% Summarize your main contributions, key findings, and their significance.
% Avoid simply repeating the abstract - provide a synthesis that reinforces your contributions.


% Bibliography
\bibliographystyle{plainnat}
\bibliography{references} % Make sure you have a references.bib file

\end{document}
